\section{Diseño de Datos}

El simulador maneja seis estructuras de datos principales: el espacio celular, la vecindad, la regla de transición, la configuración del autómata celular, la configuración del autómata celular neuronal y el historial de las evoluciones. A continuación se definen estas variables y cómo evolucionan en las diferentes capas del sistema.

La información dentro de CUDA, también llamado \textit{device}, no puede almacenar arreglos multidimensionales. Por esta razón, almacena matrices y tensores como un único arreglo unidimensional continuo y se accede a los indices de forma manual. Resulta necesario hacer un mapeo entre Python, también \textit{host}, con el \textit{device}.

\subsection{Espacio Celular}

Se define como una matriz bidimensional de dimensiones $H \times W$ con valores enteros de 8~bits, permitiendo representar hasta 256 estados distintos por celda.

\begin{itemize}
	\item \textbf{Host}: Implementado como un \texttt{numpy.ndarray} de \textit{shape} $(H, W)$ y \texttt{dtype=uint8}. Para maximizar el rendimiento, este arreglo se comparte mediante punteros, sin realizar copias, entre el \textit{widget} de renderizado, el controlador de simulación y el \textit{wrapper} de CUDA.
	
	\item \textbf{Device}: Utiliza dos búferes (\texttt{A\_d} y \texttt{B\_d}) de tipo \texttt{uint8\_t*} en la memoria global. Tienen un esquema de \textit{ping-pong buffering}, donde tras cada iteración del kernel, se intercambian los punteros de lectura y escritura, eliminando la necesidad de copiar datos entre pasos de evolución.
	
	\item \textbf{Sincronización perezosa}: El \textit{host} no se sincroniza tras cada paso del kernel. Una bandera \texttt{host\_dirty} indica si el búfer del \textit{host} está desactualizado respecto al \textit{device}; la copia \textit{device-to-host} se ejecuta únicamente cuando se invoca una operación de lectura (\texttt{getSpace}, \texttt{getCell}) o al exportar.
\end{itemize}

\subsection{Vecindad}

La vecindad define el conjunto de celdas que influyen en la transición de una celda objetivo. Se modela bajo un esquema semitotalístico organizado en grupos de desplazamientos relativos.

\begin{itemize}
	\item \textbf{Host}: Una estructura \texttt{list[list[list[int]]]} donde el nivel externo enumera los grupos, el intermedio las posiciones relativas, y el interno los componentes $[dx, dy]$ del desplazamiento.
	
	\item \textbf{Device}: La vecindad se define mediante dos arreglos: (1) \texttt{neighbourhood[]} que contiene los vectores de desplazamiento concatenados y (2) \texttt{group\_sizes[]} que almacena los tamaños de los grupos. 
	
	\item \textbf{Cotas técnicas}: Se establece un límite de $n_{\text{size}} \leq 64$ desplazamientos totales y $n_{\text{groups}} \leq 64$ grupos para cumplir con las restricciones de la memoria constante.
\end{itemize}

\subsection{Regla de Transición (CA Clásico)}

La regla es de naturaleza semitotalística: el nuevo estado depende del conteo de estados dentro de cada grupo de la vecindad. El tamaño total $r_{\text{size}}$ se determina por el producto de las combinaciones posibles de cada grupo:
\[
r_{\text{size}} = \prod_{g \in \text{grupos}} \big((\text{states}-1)\cdot |g| + 1\big)
\]
donde $|g|$ es la cardinalidad del grupo y \texttt{states} el número de estados permitidos.

El acceso a la regla se realiza mediante una suma posicional de las contribuciones de cada grupo de vecindad, calculando un índice lineal único para cada configuración de conteos.

\begin{itemize}
	\item \textbf{Host}: \texttt{numpy.ndarray} de \texttt{uint8}, tamaño $r_{\text{size}}$.
	
	\item \textbf{Device}: Debido a que $r_{\text{size}}$ puede alcanzar hasta $16,777,216$ entradas, la tabla reside en la memoria global del dispositivo (\texttt{rule\_d}).
\end{itemize}

\subsection{Configuración del Autómata Celular (CA)}

Esta estructura encapsula los parámetros estáticos que definen el comportamiento de un autómata clásico:

\begin{center}
	\begin{tabular}{c|c | p{0.42\linewidth}}
		\hline
		\textbf{Campo} & \textbf{Tipo} & \textbf{Descripción} \\
		\hline 
		\texttt{dimensions} & \texttt{(int, int)} & Dimensiones $H \times W$ del espacio. \\ \hline 
		\texttt{states} & \texttt{uint8} & Estados permitidos (1--256). \\ \hline 
		\texttt{neighbourhood} & \texttt{list[list[list[int]]]} & Definición de grupos de vecindad. \\ \hline 
		\texttt{rule} & \texttt{uint8[]} & Tabla indexada de la sección 3.3. \\ \hline 
		\texttt{boundary\_mode} & \texttt{enum} & Modos: \texttt{OPEN}, \texttt{PERIODIC}, \texttt{REFLECTIVE}, \texttt{REFLEXIVE}. \\ \hline 
		\texttt{colors} & \texttt{list[QColor]} & Mapeo de estados a colores para el renderizado. \\
	\end{tabular}
\end{center}

\subsection{Configuración del Autómata Celular Neuronal (NCA)}

El modelo neuronal sustituye la tabla de reglas por un Perceptrón Multicapa (MLP) que procesa la información de la vecindad.

\begin{itemize}
	\item \textbf{Topología del MLP}: Se define mediante \texttt{layer\_in} (entradas según vecindad), \texttt{layer\_out} (neuronas de salida), \texttt{hidden\_layers[]} (arquitectura oculta) y \texttt{activations[]} (códigos: \texttt{0=ReLU}, \texttt{1=Sigmoid}, \texttt{2=Tanh}, \texttt{3=Linear}).
	
	\item \textbf{Pesos y Sesgos}: Se organizan en un arreglo plano de tipo \texttt{float} con un vector de desplazamientos (\textit{offsets}). Esta estructura lineal favorece el acceso coalescente en la GPU y minimiza la fragmentación de memoria.
	
	\item \textbf{Formato de Intercambio}: La persistencia se realiza en archivos \texttt{.npz} (NumPy) que contienen los tensores de pesos, sesgos y los metadatos de la topología, permitiendo una carga directa en la interfaz gráfica.
\end{itemize}

\subsection{Historial de Evoluciones}

El historial gestiona el estado temporal de la simulación mediante búferes en memoria y persistencia en disco.

\begin{itemize}
	\item \textbf{Gestión de Memoria}: 
	\begin{itemize}
		\item \texttt{past\_record}: Lista de instantáneas \texttt{ndarray} para permitir el retroceso cronológico.
		\item \texttt{future\_record}: Un \texttt{collections.deque} que almacena estados precomputados asíncronamente por un \texttt{PrefetchWorker}.
	\end{itemize}
	
	\item \textbf{Grabación de Alto Rendimiento}: Se utiliza memoria \textit{pinned} (fija) reservada con \texttt{cudaHostAlloc}. Esto permite que el flujo de datos dispositivo--anfitrión sea asíncrono y no se vea penalizado por la paginación del sistema operativo.
	
	\item \textbf{Formatos de Salida}: Los datos se exportan simultáneamente en \texttt{.npy} (para re-entrenamiento en Python) y \texttt{.mat} (para análisis y validación en Wolfram Mathematica).
\end{itemize}